\clearpage\section*{Solutions EV3-T1}\label{EV3-T1-CORR}
\textbf{MCQ:} 1 A, 2 C, 3 D, 4 D, 5 C, 6 C, 7 C, 8 B, 9 B, 10 C, 11 C, 12 A. Each correct answer: 0.5 point; total 6.
\paragraph{Checks of the fluency questions.} Order of operations; sum of fractions; subtraction of exponents; multiplying coefficient 1.2; division by 0.8; product 1.25 times 0.8; solving an affine equation; vertical increment divided by horizontal; mean; angle divided by 360; complement to 1; substitution in brackets.
\subsection*{Exercise 1 -- A compound gain}
\begin{enumerate}
\item Values -2,1,4; probabilities $\dfrac{1}{4}$, $\dfrac{1}{2}$, $\dfrac{1}{4}$. \textit{Marking: 1 point. Result and basic justification expected; half a point possible for a correct but unfinished approach.}
\item $E(X)=1$; $E(X^2)=1+\dfrac{1}{2}+4=\dfrac{11}{2}$; $V(X)=\dfrac{9}{2}$. \textit{Marking: 2 points. Method or justification: 1; calculation and conclusion: 1.}
\item Mean gross gain 3: fair fee 3. Standard deviation unchanged, $\dfrac{3}{\sqrt 2}$. \textit{Marking: 2 points. Method or justification: 1; calculation and conclusion: 1.}
\end{enumerate}
\subsection*{Exercise 2 -- Peak of a model}
\begin{enumerate}
\item $f\prime=-(t+1)e^{-t}$. \textit{Marking: 1 point. Result and basic justification expected; half a point possible for a correct but unfinished approach.}
\item Strictly decreasing on $[0;+\infty[$, maximum 2 at 0. \textit{Marking: 2 points. Method or justification: 1; calculation and conclusion: 1.}
\item $y=2-t$; approximation 1.9. This is not an exact equality. \textit{Marking: 2 points. Method or justification: 1; calculation and conclusion: 1.}
\end{enumerate}
\subsection*{Exercise 3 -- The unit circle}
\begin{enumerate}
\item $A(\dfrac{\sqrt 3}{2};\dfrac{1}{2})$, $B(-\dfrac{\sqrt 3}{2};\dfrac{1}{2})$. \textit{Marking: 1 point. Result and basic justification expected; half a point possible for a correct but unfinished approach.}
\item $AB=\sqrt3$. \textit{Marking: 1 point. Result and basic justification expected; half a point possible for a correct but unfinished approach.}
\item Midpoint (0;$\dfrac{1}{2}$), AB horizontal; vertical perpendicular bisector $x=0$. \textit{Marking: 2 points. Method or justification: 1; calculation and conclusion: 1.}
\end{enumerate}
\paragraph{Calibration.} Target duration: 20--25 min for the fluency questions, 80--85 min for the exercises, 10--15 min for proofreading. These durations are editorial estimates to be measured in class. Accept any correct method; do not penalise twice a numerical error that is propagated consistently. The proposed marking scheme is not a decision of the team to harmonise marking.
